% ═══════════════════════════════════════════════════════════════
% SL-01d-ML: Stundenleistung MUSTERLÖSUNG
% ═══════════════════════════════════════════════════════════════

\documentclass[11pt, a4paper]{article}

\newif\ifswmodus
\swmodusfalse

\usepackage[utf8]{inputenc}
\usepackage[T1]{fontenc}
\usepackage[ngerman]{babel}
\usepackage{helvet}
\renewcommand{\familydefault}{\sfdefault}

\usepackage[margin=2cm]{geometry}
\usepackage{enumitem}
\usepackage{tabularx}
\usepackage{booktabs}

\usepackage{xcolor}
\usepackage{tcolorbox}
\tcbuselibrary{skins}

\usepackage{fancyhdr}
\usepackage{lastpage}
\usepackage{needspace}

\definecolor{spanisch-color}{HTML}{c41e3a}
\definecolor{grau-color}{HTML}{888888}
\definecolor{hellgrau-color}{HTML}{f5f5f5}
\definecolor{loesung-color}{HTML}{c62828}

\definecolor{sw-dark}{HTML}{000000}
\definecolor{sw-gray}{HTML}{666666}
\definecolor{sw-white}{HTML}{ffffff}

\ifswmodus
    \colorlet{spanisch}{sw-dark}
    \colorlet{grau}{sw-gray}
    \colorlet{hellgrau}{sw-white}
    \colorlet{loesung}{sw-dark}
\else
    \colorlet{spanisch}{spanisch-color}
    \colorlet{grau}{grau-color}
    \colorlet{hellgrau}{hellgrau-color}
    \colorlet{loesung}{loesung-color}
\fi

\newcommand{\fachfarbe}{spanisch}

\newcommand{\thema}{Stundenleistung: Empezamos}
\newcommand{\fach}{Spanisch}
\newcommand{\klasse}{7}
\newcommand{\maxpunkte}{20}
\newcommand{\zeit}{25 Minuten}
\newcommand{\datum}{\today}

\pagestyle{fancy}
\fancyhf{}
\renewcommand{\headrulewidth}{0.4pt}
\renewcommand{\footrulewidth}{0.4pt}

\lhead{\fach{} -- Klasse \klasse}
\chead{\textcolor{loesung}{\textbf{SL-01d MUSTERLÖSUNG}}}
\rhead{\datum}

\lfoot{\zeit}
\cfoot{}
\rfoot{Seite \thepage\ von \pageref{LastPage}}

\newcommand{\punktebox}[1]{\hfill\fbox{\small \_/#1}}
\newcommand{\luecke}[1]{\rule{#1}{0.4pt}}
\newcommand{\lsg}[1]{\textcolor{loesung}{\textbf{#1}}}

\newtcolorbox{infobox}{
    colback=hellgrau,
    colframe=\fachfarbe,
    boxrule=1pt,
    arc=0pt,
    left=8pt, right=8pt, top=8pt, bottom=8pt
}

\newenvironment{aufgabe}[1]{%
    \needspace{5\baselineskip}%
    \nopagebreak[4]%
    \vspace{10pt}
    \noindent\textbf{\textcolor{\fachfarbe}{Aufgabe #1}}
    \nopagebreak[4]%
    \vspace{5pt}
    \nopagebreak[4]%
    \begin{list}{}{\leftmargin=0pt}
    \item[]
}{%
    \end{list}
}

\newcommand{\es}[1]{\textcolor{\fachfarbe}{\textbf{#1}}}

\begin{document}

\begin{center}
    {\Large\bfseries\textcolor{\fachfarbe}{\thema}}\\[0.3em]
    {\large\textcolor{loesung}{\textbf{--- MUSTERLÖSUNG ---}}}
\end{center}

\vspace{5pt}

\noindent
\begin{tabularx}{\textwidth}{|X|X|X|}
\hline
\textbf{Name:} \lsg{Musterschüler/in} &
\textbf{Klasse:} \klasse &
\textbf{Datum:} \datum \\
\hline
\end{tabularx}

\vspace{10pt}

\begin{infobox}
\textbf{Zeit:} \zeit{} | \textbf{Punkte:} \maxpunkte{} | \textbf{Hilfsmittel:} keine
\end{infobox}

\vspace{15pt}

\begin{aufgabe}{1 -- Begrüßungen zuordnen}
Verbinde die spanische Begrüßung mit der deutschen Bedeutung. \punktebox{4}
\end{aufgabe}

\begin{center}
\begin{tabular}{rcl}
\es{¡Hola!} & \lsg{---} & Gute Nacht! \lsg{$\rightarrow$ Hola = Hallo} \\[0.8em]
\es{¡Buenos días!} & \lsg{---} & Hallo! \lsg{$\rightarrow$ Buenos días = Guten Tag} \\[0.8em]
\es{¡Buenas noches!} & \lsg{---} & Auf Wiedersehen! \lsg{$\rightarrow$ Buenas noches = Gute Nacht} \\[0.8em]
\es{¡Adiós!} & \lsg{---} & Guten Tag! \lsg{$\rightarrow$ Adiós = Auf Wiedersehen} \\
\end{tabular}
\end{center}

\vspace{15pt}

\begin{aufgabe}{2 -- Das Verb \es{ser} (sein)}
Ergänze die richtige Form von \es{ser}. \punktebox{4}
\end{aufgabe}

\noindent
a) Yo \lsg{soy} de Alemania. \hfill (ich bin)

\vspace{0.8em}

\noindent
b) Tú \lsg{eres} María. \hfill (du bist)

\vspace{0.8em}

\noindent
c) Ella \lsg{es} profesora. \hfill (sie ist)

\vspace{0.8em}

\noindent
d) Nosotros \lsg{somos} de Rostock. \hfill (wir sind)

\vspace{15pt}

\begin{aufgabe}{3 -- Artikel einsetzen}
Wähle den richtigen Artikel: \es{el}, \es{la}, \es{los} oder \es{las}. \punktebox{6}
\end{aufgabe}

\noindent
a) \lsg{el} chico (der Junge) \hfill
b) \lsg{la} chica (das Mädchen) \\[0.8em]
c) \lsg{los} libros (die Bücher) \hfill
d) \lsg{la} mesa (der Tisch) \\[0.8em]
e) \lsg{las} profesoras (die Lehrerinnen) \hfill
f) \lsg{el} cuaderno (das Heft) \\

\vspace{15pt}

\begin{aufgabe}{4 -- Dialog vervollständigen}
Ergänze den Dialog sinnvoll auf Spanisch. \punktebox{6}
\end{aufgabe}

\noindent
\textbf{A:} ¡Hola! ¿Cómo te llamas?

\vspace{0.8em}

\noindent
\textbf{B:} \lsg{Me llamo [Name]. / Soy [Name].} \hfill (2P)

\vspace{0.8em}

\noindent
\textbf{A:} ¿Qué tal?

\vspace{0.8em}

\noindent
\textbf{B:} \lsg{Bien, gracias. / Muy bien. / Regular.} \hfill (2P)

\vspace{0.8em}

\noindent
\textbf{A:} ¿De dónde eres?

\vspace{0.8em}

\noindent
\textbf{B:} \lsg{Soy de Alemania. / Soy de [Ort].} \hfill (2P)

\vspace{25pt}
\hrule
\vspace{10pt}

\noindent
\begin{tabularx}{\textwidth}{|X|c|c|c|c||c|}
\hline
\textbf{Aufgabe} & 1 & 2 & 3 & 4 & \textbf{Gesamt} \\
\hline
\textbf{max. Punkte} & 4 & 4 & 6 & 6 & \textbf{20} \\
\hline
\textbf{erreicht} & \lsg{4} & \lsg{4} & \lsg{6} & \lsg{6} & \lsg{20} \\
\hline
\end{tabularx}

\vspace{10pt}

\begin{flushright}
\textbf{Note:} \lsg{1}
\end{flushright}

\vspace{1em}

\textbf{Notenschlüssel:}

\begin{tabular}{|c|c|c|c|c|c|c|}
\hline
\textbf{Note} & 1 & 2 & 3 & 4 & 5 & 6 \\
\hline
\textbf{ab Punkte} & 19 & 16 & 12 & 8 & 4 & <4 \\
\hline
\end{tabular}

\end{document}
